\begin{proof}

\begin{enumerate}[label={$\langle1\rangle$\arabic*.},leftmargin=*]
\item Let $\eta:\widetilde S\to S$ be normalization and choose a resolution $\pi:Z\to\widetilde S$ such that the exceptional curves have normal crossings; the components of $Z$ are compact smooth surfaces. Pullback of Kähler differentials along $h=\eta\circ\pi:Z\to S$ kills the torsion of $\Omega_S^1$, because the resulting form vanishes on the dense inverse image of $S_{\rm reg}$ and the target is locally free. Thus a section $\alpha$ on the left of \eqref{cdp:eq:A.20} gives a holomorphic one-form $\alpha_Z$ with values in the unitary-flat line $h^*A$. This map on sections is injective for the same dense-open reason.

\item Write $\nabla$ for the flat unitary connection. Holomorphicity gives $\nabla\alpha_Z$ type $(2,0)$. The Hermitian pairing identifies the tensor product of the flat line with its conjugate with the trivial local system. Stokes' theorem gives, on each connected component of $Z$,
\[
 \int_Z \nabla\alpha_Z\wedge\overline{\nabla\alpha_Z}
 =\int_Z d\bigl(\alpha_Z\wedge\overline{\nabla\alpha_Z}\bigr)=0.
\]
After multiplication by the fixed dimensional sign, the integrand is the nonnegative squared norm of a $(2,0)$-form times a positive volume form. Hence $\nabla\alpha_Z=0$.

\item We prove the required local primitive assertion at a point $x\in S$. Trivialize $\mathbb A$ on a sufficiently small connected neighbourhood $V$ of $x$; complex analytic spaces are locally contractible, so such a flat trivialization exists. The section $\alpha$ has a local lift to $\Omega_S^1\otimes A$, since the latter surjects onto the torsion-free quotient. Its pullback to each reduced fibre of $h$ over $x$ is zero: the morphism of that fibre to $S$ is constant, so the induced morphism on Kähler differentials is zero.

\item For each point $\widetilde x\in\eta^{-1}(x)$, a sufficiently small regular neighbourhood of the exceptional fibre $\pi^{-1}(\widetilde x)$ in the good resolution retracts onto that fibre. This neighbourhood can be constructed by plumbing the normal disc bundles of the exceptional curves at their transverse intersection points; it is a neighbourhood of the entire fibre. A sufficiently small inverse image of a neighbourhood of $\widetilde x$ lies inside it by properness. In the smooth case one instead uses a coordinate ball. Every period of $\alpha_Z$ on the regular neighbourhood is zero, because the form is closed and its restriction to the retracting fibre is zero. Integration from a fixed point therefore produces a single-valued holomorphic primitive: a smooth primitive has differential of type $(1,0)$, so its $\bar\partial$ is zero. Normality and $\pi_*\OO_Z=\OO_{\widetilde S}$ descend that primitive to a holomorphic function $b$ on a neighbourhood of $\widetilde x$. Choose its constant so that $b(\widetilde x)=0$. Perform this independently at every point over $x$ and shrink $V$ so that all its inverse image is covered by these neighbourhoods.

\item These primitives have equal values on every fibre of $\eta$ over $V$. To see this, form the finite analytic fibre product
$W=\widetilde S\times_S\widetilde S$ with projections $p_1,p_2$. On the normalization of every positive-dimensional irreducible component of $W_{\rm red}$, the differential of $p_1^*b-p_2^*b$ is zero on a dense open subset: both differentials are pullbacks of the same local lift of $\alpha$ from $S$. The isolated singularities of $\widetilde S$ can be omitted for this comparison; holomorphicity then extends the equality. Thus the difference is constant on each such component. The same conclusion is automatic on a zero-dimensional component. Shrink $V$ once more so that every component of the finite space $W|_V$ meets the fibre over $x$; components disjoint from that fibre can be removed by properness. Our normalization $b(\widetilde x)=0$ makes all these constants zero. Hence $p_1^*b=p_2^*b$ as continuous functions, proving equality on the fibres of $\eta$.

\item A finite proper surjection is a topological quotient map. The functions $b$ therefore descend to continuous functions on $V$, holomorphic after pullback to normalization, whose differential on the regular locus is $\alpha$. On overlaps, any two such primitives differ by a locally constant function: their difference is constant on each local normalized branch, and continuity identifies the constants whenever branches meet. With the flat transition functions of $\mathbb A$, these differences define a Cech class in $H^1(S,\mathbb A)$. Changing the chosen primitives changes it by a locally constant coboundary. This defines the canonical linear map \eqref{cdp:eq:A.20}.

\item If the resulting class is zero, adjust the local primitives by locally constant sections so that they glue. Their pullback is then a global holomorphic section of the unitary-flat line on each compact component of $Z$. Its squared norm is plurisubharmonic, so the maximum principle makes the section parallel; its covariant differential is zero. Consequently $\alpha_Z=0$, and injectivity of pullback gives $\alpha=0$. This proves \eqref{cdp:eq:A.20}.

\end{enumerate}
\end{proof}
